\documentclass[final]{beamer}

\usepackage[T1]{fontenc}
\usepackage[size=custom,width=84.1,height=118.9,scale=1.18]{beamerposter}
\usetheme{gemini}
\usecolortheme{uchicago}
\usepackage{graphicx}
\usepackage{tikz}
\usetikzlibrary{arrows.meta,positioning}
\usepackage{amsmath,amssymb}
\usepackage[skins,theorems]{tcolorbox}
\usepackage{unicode-math}
\defaultfontfeatures{Scale=MatchLowercase,Ligatures=TeX}
\setmainfont{Lato}[Scale=0.90]
\setmathfont{Latin Modern Math}
\usepackage[font=small,labelfont=bf,justification=raggedright]{caption}

\newlength{\sepwidth}
\newlength{\colwidth}
\setlength{\sepwidth}{0.028\paperwidth}
\setlength{\colwidth}{0.458\paperwidth}
\newcommand{\separatorcolumn}{\begin{column}{\sepwidth}\end{column}}

\title{Basin erosion before instability}
\author{Llu\'is Torres-Hugas$^{1}$ \and Ivan Kryven$^{2}$ \and Christian Bick$^{3}$}
\institute{$^{1}$Universitat Rovira i Virgili \qquad $^{2}$Utrecht University \qquad $^{3}$Vrije Universiteit Amsterdam}
\footercontent{FisEs'26 -- Tarragona, 14--16 October 2026 \hfill lluis.torres@urv.cat}

\begin{document}

\addtobeamertemplate{headline}{}{
  \begin{tikzpicture}[remember picture,overlay]
    \node[anchor=north east,inner sep=2.2cm] at ([xshift=-0.7cm,yshift=0.4cm]current page.north east)
      {\IfFileExists{logos/logo_urv.png}{\includegraphics[height=6.6cm]{logos/logo_urv.png}}{}};
  \end{tikzpicture}
}

\begin{frame}[t]

\begin{columns}[t]
\separatorcolumn
\begin{column}{0.944\paperwidth}
\begin{alertblock}{}
\centering
{\LARGE\bfseries Finite response can erode a basin before the recovered state loses local stability.}
\end{alertblock}
\end{column}
\separatorcolumn
\end{columns}

\begin{columns}[t]
\separatorcolumn

\begin{column}{\colwidth}

\begin{block}{Finite response changes the physics}
\begin{columns}[T,totalwidth=\linewidth]
\begin{column}{0.43\linewidth}
\includegraphics[width=\linewidth]{figures/landscape.pdf}
\end{column}
\begin{column}{0.53\linewidth}
A recovery coordinate $\theta$ follows a landscape $U(\theta)$, but the responding state $x$ is not instantaneous:
\[
\dot\theta=\eta[g(\theta)-x],\qquad
\dot x=\gamma(\theta-x),
\qquad \rho=\eta/\gamma.
\]
Eliminating $x$ gives
\[
\ddot\theta+[1-\rho+\rho U''(\theta)]\dot\theta+\rho U'(\theta)=0.
\]
\textbf{Finite response turns gradient recovery into an inertial dynamics with state-dependent damping.}
\end{column}
\end{columns}

\vspace{0.8cm}
\begin{center}
\begin{tikzpicture}[node distance=1.4cm and 1.4cm,>=Latex,every node/.style={align=center}]
\node[draw=boxbg,very thick,rounded corners=2pt,inner sep=8pt] (a) {gradient\\recovery};
\node[draw=boxbg,very thick,rounded corners=2pt,inner sep=8pt,right=of a] (b) {finite\\response};
\node[draw=boxbg,very thick,rounded corners=2pt,inner sep=8pt,right=of b] (c) {negative damping\\near the barrier};
\node[draw=boxbg,very thick,rounded corners=2pt,inner sep=8pt,right=of c] (d) {basin\\erosion};
\draw[->,very thick] (a)--(b);
\draw[->,very thick] (b)--(c);
\draw[->,very thick] (c)--(d);
\end{tikzpicture}
\end{center}
\end{block}

\begin{block}{The barrier sets the onset}
For a barrier saddle $\theta_b$, the first loss of the conservative safe region occurs when the stable basin boundary changes side at the barrier. If $g'$ is maximal there,
\begin{tcolorbox}[colback=white,colframe=boxbg,boxrule=4pt,sharp corners]
\centering
{\Huge
\[
\rho^*=\frac{1}{g'(\theta_b)}
      =\frac{1}{1-U''(\theta_b)}.
\]}
\end{tcolorbox}

For the single-barrier Helmholtz landscape
\[
U(\theta)=\frac{\theta^2}{2}-\frac{\theta^3}{3},
\qquad \rho^*=\frac12.
\]

\begin{center}
\includegraphics[width=0.88\linewidth]{figures/threshold.pdf}
\end{center}

\textbf{A local property of the barrier predicts the onset of global basin loss.}
\end{block}

\begin{block}{A global robustness measure}
Let $\Omega_0$ be the conservative reference basin and $B_\rho$ the finite-response basin. Define
\[
C(\rho)=\frac{\mu(B_\rho\cap\Omega_0)}{\mu(\Omega_0)}.
\]
$C=1$ means the prescribed perturbation set is fully retained; $C<1$ measures erosion. The stable-manifold displacement links the local barrier mechanism to this global basin-volume loss.
\end{block}

\begin{block}{Take-away}
\begin{itemize}
\item Finite response creates an effective damping that can change sign across a fixed landscape.
\item Basin erosion can start before any local loss of stability of the recovered state.
\item In the barrier-controlled case, the onset is fixed exactly by the local barrier geometry.
\end{itemize}
\end{block}

\end{column}

\separatorcolumn

\begin{column}{\colwidth}

\begin{block}{The basin moves before the fixed point destabilizes}
\begin{center}
\includegraphics[width=0.98\linewidth]{figures/basin_geometry.pdf}
\end{center}

The dashed curve is the instantaneous-response boundary. The solid curve is the stable manifold of the barrier for finite response. At $\rho>\rho^*$ it enters the reference basin and removes a finite set of recoverable perturbations.

\vspace{0.5cm}
\centering
{\Large\bfseries The attractor is still locally stable. The nonlinear resilience is already smaller.}
\end{block}

\begin{block}{The mechanism is not specific to one toy landscape}
The same finite-response structure appears exactly in bistable FitzHugh--Nagumo activator--recovery dynamics. With
\[
\rho=\frac{\alpha}{b^2\beta},
\]
the bistable symmetric case gives
\[
\rho^*=\frac{1}{b}.
\]
\begin{center}
\includegraphics[width=0.88\linewidth]{figures/fhn.pdf}
\end{center}
\end{block}

\end{column}
\separatorcolumn
\end{columns}

\end{frame}
\end{document}
